\documentclass[10pt,a4paper]{amsart}
\usepackage{amsmath,amssymb,amsthm,graphicx}
% Retained minimal preamble for the Stacks Project chapter more-algebra.tex (More on Algebra).
% The chapter file pinned-original-source.tex carries only \input{preamble} on its first line, so no
% preamble text lives inside the chapter itself.  The Stacks Project uses one single shared preamble for
% every chapter; the snapshot stages it identically as sources/stacks-chow/preamble.tex and
% sources/stacks-etale-cohomology/preamble.tex.  The block below is copied mechanically, line for line,
% from sources/stacks-chow/preamble.tex lines 14-18 (the Xy-pic setup) and 154-217 (the native theorem
% declarations together with the mathematical macros).  Nothing was added, removed or redefined; the
% wrapper class and the amsmath/amssymb/amsthm/graphicx base packages are supplied by the assembler.
% Chapter: more-algebra.tex -- More on Algebra.
% For commutative diagrams we use Xy-pic
\usepackage[all]{xy}

% We use 2cell for 2-commutative diagrams.
\xyoption{2cell}

% Theorem environments.
%
\theoremstyle{plain}
\newtheorem{theorem}[subsection]{Theorem}
\newtheorem{proposition}[subsection]{Proposition}
\newtheorem{lemma}[subsection]{Lemma}

\theoremstyle{definition}
\newtheorem{definition}[subsection]{Definition}
\newtheorem{example}[subsection]{Example}
\newtheorem{exercise}[subsection]{Exercise}
\newtheorem{situation}[subsection]{Situation}

\theoremstyle{remark}
\newtheorem{remark}[subsection]{Remark}
\newtheorem{remarks}[subsection]{Remarks}

\numberwithin{equation}{subsection}

% Macros
%
\def\lim{\mathop{\mathrm{lim}}\nolimits}
\def\colim{\mathop{\mathrm{colim}}\nolimits}
\def\Spec{\mathop{\mathrm{Spec}}}
\def\Hom{\mathop{\mathrm{Hom}}\nolimits}
\def\Ext{\mathop{\mathrm{Ext}}\nolimits}
\def\SheafHom{\mathop{\mathcal{H}\!\mathit{om}}\nolimits}
\def\SheafExt{\mathop{\mathcal{E}\!\mathit{xt}}\nolimits}
\def\Sch{\mathit{Sch}}
\def\Mor{\mathop{\mathrm{Mor}}\nolimits}
\def\Ob{\mathop{\mathrm{Ob}}\nolimits}
\def\Sh{\mathop{\mathit{Sh}}\nolimits}
\def\NL{\mathop{N\!L}\nolimits}
\def\CH{\mathop{\mathrm{CH}}\nolimits}
\def\proetale{{pro\text{-}\acute{e}tale}}
\def\etale{{\acute{e}tale}}
\def\QCoh{\mathit{QCoh}}
\def\Ker{\mathop{\mathrm{Ker}}}
\def\Im{\mathop{\mathrm{Im}}}
\def\Coker{\mathop{\mathrm{Coker}}}
\def\Coim{\mathop{\mathrm{Coim}}}

% Boxtimes
%
\DeclareMathSymbol{\boxtimes}{\mathbin}{AMSa}{"02}

%
% Macros for moduli stacks/spaces
%
\def\QCohstack{\mathcal{QC}\!\mathit{oh}}
\def\Cohstack{\mathcal{C}\!\mathit{oh}}
\def\Spacesstack{\mathcal{S}\!\mathit{paces}}
\def\Quotfunctor{\mathrm{Quot}}
\def\Hilbfunctor{\mathrm{Hilb}}
\def\Curvesstack{\mathcal{C}\!\mathit{urves}}
\def\Polarizedstack{\mathcal{P}\!\mathit{olarized}}
\def\Complexesstack{\mathcal{C}\!\mathit{omplexes}}
% \Pic is the operator that assigns to X its picard group, usage \Pic(X)
% \Picardstack_{X/B} denotes the Picard stack of X over B
% \Picardfunctor_{X/B} denotes the Picard functor of X over B
\def\Pic{\mathop{\mathrm{Pic}}\nolimits}
\def\Picardstack{\mathcal{P}\!\mathit{ic}}
\def\Picardfunctor{\mathrm{Pic}}
\def\Deformationcategory{\mathcal{D}\!\mathit{ef}}

% --- appended 2026-10-12 (Chief): the project own definitions of the reference/slogan/history
% display environments, copied mechanically from the shared Stacks preamble lines 8-12
% (\usepackage{verbatim} plus the three \newenvironment lines, which the chapter preamble keeps
% verbatim from the project preamble). Nothing added, removed or redefined: these environments
% carry no mathematical content and the project itself typesets them as comments. Without these
% three lines a counted statement whose own text uses them fails with Environment slogan undefined.
\usepackage{verbatim}
\newenvironment{reference}{\comment}{\endcomment}
\newenvironment{slogan}{\comment}{\endcomment}
\newenvironment{history}{\comment}{\endcomment}
\usepackage{cleveref}
\crefname{originallemma}{original lemma}{original lemmas}
\Crefname{originallemma}{Original lemma}{Original lemmas}
\makeatletter
\newlabel{more-algebra-lemma-formally-smooth-finite-type}{{07VH}{source}{}{sourcepointer-1}{}}
\newlabel{more-algebra-lemma-formally-smooth-finite-type@cref}{{[originallemma][0][]07VH}{[1][0][]source}{}{}{}}
\makeatother
\crefname{originallemma}{original lemma}{original lemmas}
\Crefname{originallemma}{Original lemma}{Original lemmas}
\makeatletter
\newlabel{algebra-lemma-smooth-test-artinian}{{02HT}{source}{}{sourcepointer-2}{}}
\newlabel{algebra-lemma-smooth-test-artinian@cref}{{[originallemma][0][]02HT}{[1][0][]source}{}{}{}}
\makeatother
\crefname{originallemma}{original lemma}{original lemmas}
\Crefname{originallemma}{Original lemma}{Original lemmas}
\makeatletter
\newlabel{algebra-proposition-smooth-formally-smooth}{{00TN}{source}{}{sourcepointer-3}{}}
\newlabel{algebra-proposition-smooth-formally-smooth@cref}{{[originallemma][0][]00TN}{[1][0][]source}{}{}{}}
\makeatother
\crefname{originallemma}{original lemma}{original lemmas}
\Crefname{originallemma}{Original lemma}{Original lemmas}
\makeatletter
\newlabel{algebra-proposition-characterize-formally-smooth}{{031J}{source}{}{sourcepointer-4}{}}
\newlabel{algebra-proposition-characterize-formally-smooth@cref}{{[originallemma][0][]031J}{[1][0][]source}{}{}{}}
\makeatother
\crefname{originallemma}{original lemma}{original lemmas}
\Crefname{originallemma}{Original lemma}{Original lemmas}
\makeatletter
\newlabel{algebra-lemma-NL-localize-bottom}{{07BS}{source}{}{sourcepointer-5}{}}
\newlabel{algebra-lemma-NL-localize-bottom@cref}{{[originallemma][0][]07BS}{[1][0][]source}{}{}{}}
\makeatother
\crefname{originallemma}{original lemma}{original lemmas}
\Crefname{originallemma}{Original lemma}{Original lemmas}
\makeatletter
\newlabel{algebra-lemma-localize-NL}{{00S7}{source}{}{sourcepointer-6}{}}
\newlabel{algebra-lemma-localize-NL@cref}{{[originallemma][0][]00S7}{[1][0][]source}{}{}{}}
\makeatother
\crefname{originallemma}{original lemma}{original lemmas}
\Crefname{originallemma}{Original lemma}{Original lemmas}
\makeatletter
\newlabel{lemma-formally-smooth}{{07EC}{source}{}{sourcepointer-7}{}}
\newlabel{lemma-formally-smooth@cref}{{[originallemma][0][]07EC}{[1][0][]source}{}{}{}}
\makeatother
\title{Smoothness at a prime of a finite type algebra is formal smoothness of the localisation}
\author{The Stacks Project Authors}
\date{stacks-project:more-algebra.tex:local-snapshot-20261012; GNU FDL 1.2 or later}
\begin{document}
\maketitle
\setcounter{subsection}{1}
\noindent\textit{Editorial note.} All delivered mathematical text is the author's own, copied from the pinned Stacks Project chapter file more-algebra.tex; no statement or proof step was rewritten, reordered or removed.\par\medskip
\noindent\textit{Editorial note.} The chapter file pinned-original-source.tex carries only \textbackslash{}input\{preamble\} on its first line, so no preamble text lives inside the chapter. The retained file sources/stacks-more-algebra/retained-preamble.tex is a mechanical, line-for-line copy of the project's single shared preamble as staged in the snapshot (Xy-pic setup plus theorem declarations and mathematical macros). Nothing was added, removed or redefined, and the excerpt serves only as preamble, never as a span coordinate.\par\medskip
\noindent\textit{Editorial note.} Both counted spans start at the \textbackslash{}begin line of their environment, so the statement and the author's proof are delivered whole; surrounding prose of the section is deliberately outside the counted spans and is carried as standard prerequisites instead.\par\medskip
\noindent\textit{Editorial note.} The citation fields (source\_id, source\_version, authors, source\_url, license.name, source\_title) are identical to those of every other card this lane has produced from this chapter, so the batch selector treats them as one source.\par\medskip
\noindent\textit{Allowed standard prerequisites (editorial).}\begin{itemize}
\item The Artinian test criterion for smoothness of a finite type algebra at a prime. Algebra, Lemma \textbackslash{}ref\{algebra-lemma-smooth-test-artinian\}, cited at proof lines 9667-9668. Gives the implication from formal smoothness to smoothness.
\item A smooth algebra map is formally smooth. Algebra, Proposition \textbackslash{}ref\{algebra-proposition-smooth-formally-smooth\}, cited at proof lines 9670-9671. Provides formal smoothness of the localised algebra.
\item Formal smoothness is preserved by localisation, via the characterisation through the naive cotangent complex. Algebra, Proposition \textbackslash{}ref\{algebra-proposition-characterize-formally-smooth\}, Algebra, Lemma \textbackslash{}ref\{algebra-lemma-NL-localize-bottom\}, and Algebra, Lemma \textbackslash{}ref\{algebra-lemma-localize-NL\}, cited at proof lines 9673-9680. Justifies passing to the localisations at the two primes.
\item Formal smoothness upgrades to the adic topology. Lemma \textbackslash{}ref\{lemma-formally-smooth\}, cited at proof lines 9681-9683. Turns algebraic formal smoothness into the adic statement of the lemma.
\end{itemize}\medskip
\ifdefined\hypertarget\hypertarget{sourcepointer-1}{}\fi
\noindent\textit{Source reference.} Original lemma 07VH: focused-seeds/AGP/sources/stacks-snapshot-20261011/tags.txt, source lines 10007--10007. The lemma's own canonical tag, quoted so that readers can resolve the label inside the delivered text; it points at the counted statement itself and carries no proof content.\par\medskip
\ifdefined\hypertarget\hypertarget{sourcepointer-2}{}\fi
\noindent\textit{Source reference.} Original lemma 02HT: focused-seeds/AGP/sources/stacks-snapshot-20261011/tags.txt, source lines 3276--3276. Navigation pointer to a statement the counted proof refers to but does not contain; nothing load-bearing is delegated to it.\par\medskip
\ifdefined\hypertarget\hypertarget{sourcepointer-3}{}\fi
\noindent\textit{Source reference.} Original lemma 00TN: focused-seeds/AGP/sources/stacks-snapshot-20261011/tags.txt, source lines 1081--1081. Navigation pointer to a statement the counted proof refers to but does not contain; nothing load-bearing is delegated to it.\par\medskip
\ifdefined\hypertarget\hypertarget{sourcepointer-4}{}\fi
\noindent\textit{Source reference.} Original lemma 031J: focused-seeds/AGP/sources/stacks-snapshot-20261011/tags.txt, source lines 3986--3986. Navigation pointer to a statement the counted proof refers to but does not contain; nothing load-bearing is delegated to it.\par\medskip
\ifdefined\hypertarget\hypertarget{sourcepointer-5}{}\fi
\noindent\textit{Source reference.} Original lemma 07BS: focused-seeds/AGP/sources/stacks-snapshot-20261011/tags.txt, source lines 9352--9352. Navigation pointer to a statement the counted proof refers to but does not contain; nothing load-bearing is delegated to it.\par\medskip
\ifdefined\hypertarget\hypertarget{sourcepointer-6}{}\fi
\noindent\textit{Source reference.} Original lemma 00S7: focused-seeds/AGP/sources/stacks-snapshot-20261011/tags.txt, source lines 1029--1029. Navigation pointer to a statement the counted proof refers to but does not contain; nothing load-bearing is delegated to it.\par\medskip
\ifdefined\hypertarget\hypertarget{sourcepointer-7}{}\fi
\noindent\textit{Source reference.} Original lemma 07EC: focused-seeds/AGP/sources/stacks-snapshot-20261011/tags.txt, source lines 9442--9442. Navigation pointer to a statement the counted proof refers to but does not contain; nothing load-bearing is delegated to it.\par\medskip
\begin{lemma}
\label{lemma-formally-smooth-finite-type}
Let $A \to B$ be a finite type ring map with $A$ Noetherian.
Let $\mathfrak q \subset B$ be a prime ideal lying over
$\mathfrak p \subset A$. The following are equivalent
\begin{enumerate}
\item $A \to B$ is smooth at $\mathfrak q$, and
\item $A_\mathfrak p \to B_\mathfrak q$ is formally smooth in
the $\mathfrak q$-adic topology.
\end{enumerate}
\end{lemma}

\begin{proof}
The implication (2) $\Rightarrow$ (1) follows from
Algebra, Lemma \ref{algebra-lemma-smooth-test-artinian}.
Conversely, if $A \to B$ is smooth at $\mathfrak q$, then
$A \to B_g$ is smooth for some $g \in B$, $g \not \in \mathfrak q$.
Then $A \to B_g$ is formally smooth by
Algebra, Proposition \ref{algebra-proposition-smooth-formally-smooth}.
Hence $A_\mathfrak p \to B_\mathfrak q$ is formally smooth
as localization preserves formal smoothness (for example by
the criterion of Algebra, Proposition
\ref{algebra-proposition-characterize-formally-smooth}
and the fact that the cotangent complex behaves well with respect to
localization, see
Algebra, Lemmas \ref{algebra-lemma-NL-localize-bottom} and
\ref{algebra-lemma-localize-NL}).
Finally, Lemma \ref{lemma-formally-smooth} implies that
$A_\mathfrak p \to B_\mathfrak q$ is formally smooth in the
$\mathfrak q$-adic topology.
\end{proof}

\end{document}
